\section{Algebraization and the affine counterexample}
\label{algebra:section}

We now pass from the analytic germ to an affine variety.  Approximating
the defining equations does not by itself retain their relations.
We approximate the equations and a finite relation matrix together;
their exact zero product, together with sufficiently close approximation,
will recover the same complete local ring.  Tangent freeness then descends
from the continuous derivations of that ring.  This follows the
equation-and-relation determinacy method of Hironaka
\cite[\S3, Theorem~3.3]{Hironaka1969}, used in Artin's
isolated-singularity algebraization
\cite[Theorems~3.8--3.9]{Artin1969}.  We prove the finite-determinacy
statement needed here, including fixed estimates that avoid a
complete-intersection hypothesis.

\subsection{Finite determinacy of equations with relations}

Put
\[
 P=\CC[[x_1,\ldots,x_e]],\qquad
 \mathfrak n=(x_1,\ldots,x_e),\qquad
 I=(f_1,\ldots,f_s),\qquad R=P/I,
\]
where $f=(f_1,\ldots,f_s)\in\mathfrak nP^s$ is a row.  Write
$\mathfrak m_R=\mathfrak nR$, and let an overbar denote reduction modulo $I$.
Choose an $s\times t$ matrix $D$ whose columns generate all the relations
on $f$.  Thus
\begin{equation}
\label{algebra:ideal-presentation}
 P^t\xrightarrow{D}P^s\xrightarrow{f}I\longrightarrow0
\end{equation}
is exact.  Define
\[
 \psi:R^s\longrightarrow R^t,\qquad v\longmapsto v\bar D,
 \qquad K=\ker\psi,
 \qquad
 J=\sum_{j=1}^e R\,\overline{\partial_j f}.
\]
Differentiating $fD=0$ and reducing modulo $I$ shows that $J\subset K$.
Right-exact reduction of \eqref{algebra:ideal-presentation} gives a
presentation of $I/I^2$; its dual therefore identifies
\begin{equation}
\label{algebra:conormal-dual}
 K=\Hom_R(I/I^2,R),\qquad T^1=K/J.
\end{equation}
In particular, no exactness on the left after reduction has been assumed.

\begin{lemma}[Finite determinacy with a relation matrix]
\label{algebra:finite-determinacy}
Suppose that $T^1$ in \eqref{algebra:conormal-dual} has finite length.
There exist integers $N,n_0$ depending only on $f,D$ such that the following
holds.  If a row $g\in P^s$ and an $s\times t$ matrix $B$ satisfy
\[
 gB=0,\qquad
 g-f\in\mathfrak n^{n_0}P^s,\qquad
 B-D\in\mathfrak n^NP^{s\times t},
\]
then there are a continuous $\CC$-algebra automorphism $\Phi$ of $P$ and
an invertible matrix $U\in\operatorname{GL}_s(P)$ such that
\[
 f=\Phi(g)U.
\]
One can take $\Phi(x_j)\equiv x_j\pmod{\mathfrak n^2}$ and
$U\equiv1\pmod{\mathfrak n}$.  Consequently $P/(g)\simeq P/(f)$ as
complete local $\CC$-algebras.
\end{lemma}

\begin{proof}
We improve the equations one order at a time.  The exact relation
$gB=0$ makes the equation error, reduced modulo $I$, nearly belong to
$K$.  Artin--Rees lets us correct it into $K$; the finite length of
$K/J$ then expresses the sufficiently high-order corrected class by
coordinate derivatives.  A change of generators removes the remaining
$I$-valued error.  We arrange these corrections so that the equation
precision increases at every step while the relation precision stays
fixed.  The resulting coordinate and generator changes converge in $P$.

We first choose the fixed bounds. For a map
$u:M\to M'$ of finite modules over a Noetherian local ring with maximal
ideal $\mathfrak a$, the Artin--Rees lemma~\cite[Tag~00IN]{StacksCompletion}, applied to
$\operatorname{im}u\subset M'$ gives a constant $a\ge0$ such that
\[
 \operatorname{im}u\cap\mathfrak a^rM'
 \subset u(\mathfrak a^{r-a}M)\qquad(r\ge a).
\]
Choose nonnegative constants $\alpha,\beta,d$ so that
\begin{align}
 \operatorname{im}\psi\cap\mathfrak m_R^rR^t
 &\subset\psi(\mathfrak m_R^{r-\alpha}R^s),
 \label{algebra:ar-image}\\
 K\cap\mathfrak m_R^rR^s
 &\subset\mathfrak m_R^{r-\beta}K,
 \label{algebra:ar-kernel}\\
 I^{\oplus s}\cap\mathfrak n^rP^s
 &\subset\{fC:C\in\mathfrak n^{r-d}P^{s\times s}\}
 \label{algebra:ar-ideal}
\end{align}
whenever the corresponding exponents are nonnegative.  In
\eqref{algebra:ar-ideal}, Artin--Rees is applied to the map
$P^{s\times s}\to P^s$, $C\mapsto fC$, with image $I^{\oplus s}$; it is the
filtration induced from $P^s$ that is relevant.

Choose $\ell\ge0$ with $\mathfrak m_R^\ell K\subset J$ and set
\begin{equation}
\label{algebra:constants}
 c_0=\beta+\ell,\qquad c=c_0+d,\qquad
 N\ge\alpha+1,\qquad
 n_0\ge\max\{c+N,c+2,2c+1\}.
\end{equation}
We first describe one improvement step.  Suppose $n\ge n_0$ and
\[
 h=g-f\in\mathfrak n^nP^s,
 \qquad B-D\in\mathfrak n^NP^{s\times t},
 \qquad gB=0.
\]
The exact relation, reduced modulo the fixed ideal $I$, gives
\begin{equation}
\label{algebra:high-order-relation}
 \bar h\bar D
 =\bar g\bar D
 =-\bar g(\bar B-\bar D)
 =-\bar h(\bar B-\bar D)
 \in\mathfrak m_R^{n+N}R^t.
\end{equation}
The factor $\bar h$ on the right is what provides order $n+N$.
By \eqref{algebra:ar-image}, there is
\[
 \bar\epsilon\in\mathfrak m_R^{n+N-\alpha}R^s
 \subset\mathfrak m_R^{n+1}R^s,
 \qquad \psi(\bar\epsilon)=\psi(\bar h).
\]
Hence
\[
 \kappa=\bar h-\bar\epsilon
 \in K\cap\mathfrak m_R^nR^s
 \subset\mathfrak m_R^{n-\beta}K
 \subset\mathfrak m_R^{n-c_0}J.
\]
Here the last inclusion uses $\mathfrak m_R^\ell K\subset J$,
with the remaining factor of order $n-\beta-\ell=n-c_0$.
Lift coefficients expressing $\kappa$ in the displayed generators of
$J$ to elements $a_j\in\mathfrak n^{n-c_0}$, and lift
$\bar\epsilon$ to $\epsilon\in\mathfrak n^{n+1}P^s$.  Then
\[
 w=h-\sum_{j=1}^e a_j\partial_jf-\epsilon
 \in I^{\oplus s}\cap\mathfrak n^{n-c_0}P^s.
\]
The order of this ideal-valued remainder need only be $n-c_0$, since
the lifted derivative combination may have lower order than $h$:
a derivative of a defining equation can be a unit.
Using \eqref{algebra:ar-ideal}, write $w=fC$ with
$C\in\mathfrak n^{n-c}P^{s\times s}$.  We have proved
\begin{equation}
\label{algebra:linear-correction}
 h=\sum_{j=1}^e a_j\partial_jf+fC+\epsilon,
 \qquad
 a_j,C\in\mathfrak n^{n-c},
 \qquad \epsilon\in\mathfrak n^{n+1}P^s.
\end{equation}
Thus the loss $d$ for lifting the ideal-valued remainder has been included
in the fixed constant $c$.

Put $q=n-c$, and define
\begin{equation}
\label{algebra:one-step}
 \phi(x_j)=x_j-a_j,\qquad M=(1+C)^{-1},\qquad
 g_{\mathrm{new}}=\phi(g)M,\qquad
 B_{\mathrm{new}}=(1+C)\phi(B).
\end{equation}
Our bounds ensure
\[
 q\ge\max\{N,2\},\qquad 2q\ge n+1,
 \qquad n+q-1\ge n+1.
\]
In particular, $\phi$ is a formal automorphism with identity linear part,
and $M$ is invertible.  Formal Taylor expansion gives
\[
 \phi(f)-f+\sum_j a_j\partial_jf\in\mathfrak n^{2q}P^s,
 \qquad
 \phi(h)-h\in\mathfrak n^{n+q-1}P^s.
\]
Combining these estimates with \eqref{algebra:linear-correction} gives
\[
 \phi(g)\equiv f(1+C)\pmod{\mathfrak n^{n+1}P^s},
 \qquad
 g_{\mathrm{new}}-f\in\mathfrak n^{n+1}P^s.
\]
The relations in \eqref{algebra:one-step} transform exactly:
\[
 g_{\mathrm{new}}B_{\mathrm{new}}
 =\phi(g)(1+C)^{-1}(1+C)\phi(B)=0.
\]
Moreover, substitution by $\phi$ changes any power series by an element
of $\mathfrak n^q$, and $C$ has order at least $q$.  Therefore
\[
 B_{\mathrm{new}}-D
 =C\phi(B)+(\phi(B)-B)+(B-D)
 \in\mathfrak n^NP^{s\times t}.
\]
The fixed relation precision $N$ is preserved, while the equation
precision increases from $n$ to $n+1$.

Starting with $(g_{n_0},B_{n_0})=(g,B)$, repeat this step for
$n=n_0,n_0+1,\ldots$.  Write $\phi_n,M_n$ for its corrections.
For completeness, the cumulative coordinate and generator changes obey
\[
 \Phi_{n_0}=\id,\quad U_{n_0}=1,\qquad
 \Phi_{n+1}=\phi_n\circ\Phi_n,\qquad
 U_{n+1}=\phi_n(U_n)M_n,
\]
so that $g_n=\Phi_n(g)U_n$ at every stage.  Since
\[
 (\phi_n-\id)(P)\subset\mathfrak n^{n-c},
 \qquad M_n-1\in\mathfrak n^{n-c}P^{s\times s},
\]
the consecutive images $\Phi_n(x_j)$ and the consecutive matrices
$U_n$ differ in orders tending to infinity.  They converge in the
complete ring $P$ to $\Phi(x_j)$ and $U$.  The resulting substitution
$\Phi$ has identity linear part and hence is an automorphism, while
$U\equiv1\pmod{\mathfrak n}$ is invertible.  Taking limits in
$g_n=\Phi_n(g)U_n$ yields $f=\Phi(g)U$, as required.
\end{proof}

The original matrix $D$ must generate all the relations, in order that
\eqref{algebra:conormal-dual} hold.  No such assertion is needed for the
approximating matrix $B$: its exact zero product and its closeness to
$D$ are sufficient.  Thus Lemma~\ref{algebra:finite-determinacy} concerns
solutions of the finite system $gB=0$, not arbitrary truncations of a
list of defining equations.

\subsection{The algebraic neighborhood}

\begin{proposition}[Algebraization of an isolated analytic germ]
\label{prop:algebraization}
Let $(X,x)$ be a normal complex analytic germ of dimension $d\ge1$,
smooth away from $x$, and suppose that its holomorphic tangent module
is free of rank $d$.  There exist a finitely generated normal integral
$\CC$-algebra $A$ and a maximal ideal $\mathfrak p\subset A$ such that
\[
 \dim A=d,\qquad A/\mathfrak p=\CC,\qquad
 \widehat{A_{\mathfrak p}}\simeq\widehat{\OO_{X,x}},\qquad
 \Der_{\CC}(A)\simeq A^d.
\]
The affine variety $\Spec A$ can be chosen smooth away from
$\mathfrak p$.  If $x$ is singular, $A_{\mathfrak p}$ is nonregular.
\end{proposition}

\begin{proof}
\emph{The completed analytic presentation.}
Embed the germ in a smooth analytic germ and write
\[
 S=\OO_{X,x}=\CC\{x_1,\ldots,x_e\}/I_{\mathrm{an}}.
\]
Choose a finite row $f$ generating $I_{\mathrm{an}}$ and a finite matrix
$D$ generating all its relations over the convergent power-series ring.
Completion is exact on finite modules, so they give the presentation
\eqref{algebra:ideal-presentation} over $P=\CC[[x_1,\ldots,x_e]]$,
with $R=\widehat S=P/(f)$.  Complex analytic local rings are excellent
\cite[Proposition~2.1(1)]{Cutkosky2017};
thus the completion of the normal local ring $S$ is a normal local
domain.  Also
\begin{equation}
\label{algebra:completion-invariants}
 \dim R=\dim S=d,
 \qquad R\text{ is regular if and only if }S\text{ is regular}.
\end{equation}
We use here the usual exactness and permanence properties of completion
\cite{StacksCompletion}.

The convergent version of \eqref{algebra:conormal-dual} is the conormal
dual of the analytic embedding.  On the smooth locus the tangent-to-normal
map is surjective: a smooth analytic subspace of a complex manifold is
locally a regular embedding.  Its cokernel is consequently a coherent
module supported at the isolated point $x$, and has finite length.
Flat completion identifies this cokernel with the module $T^1=K/J$
above.  Lemma~\ref{algebra:finite-determinacy} therefore applies to $f,D$.

We also record precisely which differentials are being completed.  The
module of holomorphic differentials of $S$ has the finite Jacobian
presentation with generators $dx_j$ and relations $df_i$.  Its dual is
the holomorphic tangent module.  Flat completion of this kernel yields
\begin{equation}
\label{algebra:continuous-tangent}
 \Der^{\mathrm{cont}}_{\CC}(R,R)
 =\ker\left(R^e\longrightarrow R^s,
       (v_j)\longmapsto\sum_{j=1}^e v_j\overline{\partial_jf}\right)
 \simeq R^d.
\end{equation}
Here ``continuous'' refers to the maximal-ideal topology.  A continuous
derivation of a power-series quotient is determined by its values on
the $x_j$, and the displayed relations are precisely the conditions for
it to descend from $P$.  All completed differential modules used below
are these finite continuous differential modules.

\emph{Approximation of equations and relations.}
Let
\[
 H=\bigl(\CC[x_1,\ldots,x_e]_{(x_1,\ldots,x_e)}\bigr)^h
\]
be the henselization, whose completion is $P$.  This is precisely a
henselization of a finite-type algebra over the field $\CC$ at a prime,
as covered by Artin's Theorem~1.10.  Apply that approximation theorem
over $H$ to the polynomial system in the $s+st$ entries of a row $G$
and a matrix $B$,
\begin{equation}
\label{algebra:artin-system}
 \sum_{i=1}^s G_i B_{i\nu}=0\qquad(1\le\nu\le t),
\end{equation}
with formal solution $(G,B)=(f,D)$.  It gives a row $g$ and a matrix $B$
with entries in $H$, satisfying $gB=0$ exactly and agreeing with $f,D$
to any prescribed finite order
\cite[Theorem~1.10]{Artin1969}; see also~\cite{StacksApproximation}.  Choose that
order at least $\max\{n_0,N\}$, with $n_0,N$ from
Lemma~\ref{algebra:finite-determinacy}.  The lemma gives
\begin{equation}
\label{algebra:formal-equivalence}
 P/(g)\simeq R.
\end{equation}

The henselization is the filtered colimit of pointed \'etale
neighborhoods of the origin, with the same residue field
\cite{StacksHenselization}.  Thus the finitely many entries of $g,B$
are represented on a single such neighborhood.  After passing to a
later neighborhood, the finitely many identities
\eqref{algebra:artin-system} hold there as well.  If this stage is
expressed over the local polynomial ring, finite presentation and
clearing finitely many denominators spread it to a neighborhood of
the origin in affine space; shrinking preserves the marked point
and makes the morphism \'etale.  After an affine shrinking the
identities hold in its coordinate ring.  We obtain an affine
\'etale neighborhood $\Spec E\to\mathbb A^e_{\CC}$ with a marked point
$\mathfrak q$ over the origin, residue field $\CC$, and a row $g$ of
elements of $E$ vanishing at $\mathfrak q$.  Its completed ambient
local ring is $P$.  Put
\[
 A'=E/(g),\qquad \mathfrak p'=\mathfrak q/(g),\qquad
 T=A'_{\mathfrak p'}.
\]
Then $A'$ is of finite type over $\CC$, and exactness of completion
together with \eqref{algebra:formal-equivalence} gives
\begin{equation}
\label{algebra:completed-local-ring}
 \widehat T=P/(g)\simeq R.
\end{equation}
In particular, this construction preserves the complete local ring,
including its dimension and its normality.

\emph{Descent of tangent freeness and local ring properties.}
Because the ambient neighborhood is \'etale over $\mathbb A^e_{\CC}$,
its differentials have basis $dx_1,\ldots,dx_e$ near the marked point.
Completing the finite presentation of $\Omega^1_{T/\CC}$ therefore
gives the continuous Jacobian presentation of $P/(g)$.  For a finitely
presented module, flat base change commutes with its dual: this follows
by expressing the dual as the kernel of a map between finite free
modules.  The continuous chain rule under
\eqref{algebra:completed-local-ring} and
\eqref{algebra:continuous-tangent} consequently give
\begin{equation}
\label{algebra:tangent-descent}
 \Der_{\CC}(T)\otimes_T\widehat T
 \simeq\Der^{\mathrm{cont}}_{\CC}(\widehat T,\widehat T)
 \simeq\widehat T^{d}.
\end{equation}
This is an equality obtained from finite presentations; it does not
identify ordinary algebraic differentials of a power-series ring with
its continuous differentials.

Let $M=\Der_{\CC}(T)$.  By \eqref{algebra:tangent-descent}, its residue
vector space has dimension $d$.  Lift a basis to $d$ elements of $M$.
Nakayama's lemma gives a surjection $T^d\to M$, which after completion
is a surjection between free modules of rank $d$ over the local ring
$\widehat T$, and hence an isomorphism.  Faithful flatness of completion
annihilates its kernel, proving $M\simeq T^d$.

The ring $T$ injects into its domain completion, so it is a domain.
Normality also descends.  Indeed, if $a/b\in\operatorname{Frac}(T)$
is integral over $T$, it is integral over the normal domain
$\widehat T$ and hence belongs to $\widehat T$.  Thus
$a\in b\widehat T\cap T=bT$, where the last equality follows from
faithful flatness, and $a/b\in T$.  Finally, completion preserves both
Krull dimension and the dimension of the maximal ideal modulo its
square.  Equations \eqref{algebra:completion-invariants} and
\eqref{algebra:completed-local-ring} show that $T$ has dimension $d$
and is nonregular whenever $S$ is nonregular.

The isolated-singularity property can also be retained.  Let
$\mathfrak j_{\mathrm{an}}$ be the $d$th Fitting ideal of the
holomorphic differential module of $S$.  The analytic Jacobian
criterion identifies its zero set with the nonsmooth locus, so some
power of the maximal ideal of $S$ is contained in
$\mathfrak j_{\mathrm{an}}$; in the smooth case this ideal is the unit
ideal.  Fitting ideals commute with base change.  The completed
Jacobian presentations and \eqref{algebra:completed-local-ring} thus
imply, by contraction from the faithfully flat completion, that
\begin{equation}
\label{algebra:isolated-descent}
 \mathfrak m_T^b\subset
 \operatorname{Fitt}_d(\Omega^1_{T/\CC})
 \quad\text{for some }b\ge0.
\end{equation}
On the reduced pure-$d$-dimensional neighborhood constructed below,
the algebraic Jacobian criterion interprets this containment as
nonsmoothness supported at the marked point.

\emph{An integral normal affine neighborhood with a free tangent module.}
The finite module $\Der_{\CC}(A')$ localizes at $\mathfrak p'$ to
$\Der_{\CC}(T)$, because $\Omega^1_{A'/\CC}$ is finitely presented.
Represent the basis just obtained after inverting an element outside
$\mathfrak p'$.  It defines a map from a free module of rank $d$ to
the derivation module whose kernel and cokernel vanish at
$\mathfrak p'$.  Their supports are closed, so a further principal
shrinking makes this map an isomorphism.

Since $T$ is a domain, exactly one irreducible component of $\Spec A'$
passes through $\mathfrak p'$, and the support of the nilradical avoids
$\mathfrak p'$.  Remove the other components and that support.  The
remaining neighborhood is reduced and irreducible.  Its normal locus
is open, since schemes of finite type over $\CC$ are excellent, and
contains $\mathfrak p'$ by the normality of $T$; restrict to this
normal locus as well.  This irreducible neighborhood has dimension
$d$, since its marked point is closed and its local ring has dimension
$d$.  Its nonsmooth locus is closed and, by
\eqref{algebra:isolated-descent}, has no positive-dimensional component
through $\mathfrak p'$.  Remove any components of that locus that do
not contain $\mathfrak p'$.

All the restrictions just made are open neighborhoods of
$\mathfrak p'$.  Choose a principal affine neighborhood contained in
their intersection, and write it as $\Spec A$, with marked maximal
ideal $\mathfrak p$.  The algebra $A$ is finitely generated, normal
and integral, its dimension is $d$, and
\[
 A_{\mathfrak p}=T,\qquad
 \Der_{\CC}(A)\simeq A^d.
\]
The marked residue field and completed local ring are unchanged, and
the variety is smooth away from that point.  This proves all the
assertions.
\end{proof}

\begin{proof}[Proof of Theorem~\ref{thm:main}]
Proposition~\ref{prop:analytic-counterexample} supplies a singular
normal complex analytic surface germ, smooth away from its marked
point, whose holomorphic tangent module is free of rank two.
Apply Proposition~\ref{prop:algebraization} and fix any of the
approximations and affine neighborhoods furnished by its proof.  Its
coordinate ring $A$ is a finitely generated normal integral
$\CC$-algebra of dimension two, with
\[
 \Der_{\CC}(A)\simeq A^2,
 \qquad A_{\mathfrak p}\text{ nonregular}.
\]
Thus the derivation module is free, and in particular projective,
whereas $A$ is not regular.  This is the required counterexample over
$\CC$.
\end{proof}

\subsection{The smooth locus and derivation blowups}

Let $X=\Spec A$ be the affine surface just constructed, and let
$U=X_{\mathrm{reg}}$ be its smooth locus.  Recall that the logarithmic
Kodaira dimension of a smooth quasiprojective surface is
\[
 \overline\kappa(U)=\kappa(K_{\overline U}+D),
\]
where $\overline U$ is a smooth projective compactification and its
reduced boundary $D=\overline U\setminus U$ is a simple normal-crossing
divisor.  This value is independent of the compactification.

\begin{corollary}\label{algebra:log-kodaira}
The smooth locus $U$ of the affine surface $X=\Spec A$ constructed in
Theorem~\ref{thm:main} has $\overline\kappa(U)=2$.
\end{corollary}

\begin{proof}
The free tangent sheaf on $X$ restricts to a trivial tangent bundle on
$U$.  If $\overline\kappa(U)\le1$, the theorem of Biswas, Gurjar and
Kolte~\cite[Theorem~0.1(1)]{BiswasGurjarKolte2014} would make the normal
affine surface $X$ smooth, contrary to Theorem~\ref{thm:main}.
The only remaining value of logarithmic Kodaira dimension for a surface
is two.
\end{proof}

For a finite $A$-module $M$ of generic rank $r$, viewed as a coherent
sheaf on $X$, its blowup is universal among proper birational morphisms
from integral varieties to $X$ for which the pullback of $M$ modulo
torsion is locally free of rank $r$.  Every such morphism factors uniquely
through the blowup.  We use Villamayor's description, as recalled in
\cite[\S1.6, Theorem~1.17]{BarajasChavezRomano2026}.

\begin{corollary}\label{algebra:derivation-blowups}
Starting with the singular surface $X=\Spec A$ of
Theorem~\ref{thm:main}, every iteration of blowing up at the derivation
module is isomorphic to $X$.  The same holds if each blowup is followed
by normalization.  Neither procedure resolves its singular point.
\end{corollary}

\begin{proof}
Since $\Der_{\CC}(A)\simeq A^2$, the identity of $X$ already has the
universal property just stated.  Its blowup at the derivation module
is therefore $X$ itself, and the operation remains unchanged at every
iteration.  Normalization also leaves $X$ unchanged because $A$ is
normal.  The relation with the Lipman--Zariski conjecture was noted in
\cite[Introduction, after Question~1]{BarajasChavezRomano2026}; the
surface above gives a negative answer to that question in characteristic
zero.
\end{proof}
